% ============================================================================
%  Chapter 1 — Introduction
%  Thesis: Generative AI, Privacy & Personal Data — TADJ Mohamed Walid, ESI SBA
%
%  DUAL-MODE FILE (standalone OR \input via \ThesisMaster). Sets \label{ch:intro}.
%  Drafted from ch1_intro_notes.md AFTER the body, so promises match deliveries.
%  GOVERNING RULE: every contribution/finding claimed here is delivered by a
%  committed chapter (Ch2-7), cited by \ref. Non-claims are explicit. The
%  methodological-headline generalization is framed as the thesis's central
%  ARGUMENT (it is argued in Ch5/Ch7), not a proven universal. Honest framings
%  ("convergence not invention", "~40% not elimination", "demonstration not
%  certification") are load-bearing and appear verbatim in substance.
% ============================================================================
\ifdefined\ThesisMaster\else
  \documentclass[11pt,a4paper]{report}
  \usepackage[utf8]{inputenc}
  \usepackage[T1]{fontenc}
  \usepackage{lmodern}
  \usepackage[margin=2.5cm]{geometry}
  \usepackage{booktabs}
  \usepackage{tabularx}
  \usepackage{array}
  \newcolumntype{L}{>{\raggedright\arraybackslash}X}
  \usepackage{amsmath,amssymb}
  \usepackage{enumitem}
  \usepackage{microtype}
  \usepackage[hidelinks]{hyperref}
  \begin{document}
\fi

\chapter{Introduction}
\label{ch:intro}

\section{Context and motivation}
\label{sec:intro-context}

Organisations hold years of sensitive internal data --- human-resources files, medical reports,
internal memoranda, system configurations --- and increasingly want it conversationally
accessible. Retrieval-augmented generation (RAG) makes that possible: documents are embedded
into a vector store, retrieved on demand, and used to ground a large language model's answers.
The appeal is real, and so is a privacy problem that hides in plain sight. The retriever does
not know who is allowed to see what, and neither does the model; both are engineered to be
\emph{helpful}, not \emph{minimally informative}. When an employee asks for a colleague's salary,
the pipeline retrieves the record and the model relays it; when an outside attacker hides an
instruction inside a document the retriever later fetches, the same pipeline can be steered into
disclosure \cite{greshake2023ipi}. Two separately documented phenomena --- language models
memorising and reproducing training data \cite{carlini2021extracting}, and adversarial prompts
steering model behaviour --- compound in RAG into a third, because the retriever itself becomes an
attacker-influenceable channel into the model's context.

The consequences are governed by data-protection law (the GDPR in particular), under which
unauthorised disclosure of personal data is a reportable harm, and they are made harder by the
nature of the worst leak. The most damaging disclosure in this setting is not a single record
read out verbatim but \emph{re-identification by linkage}: a person reconstructed from
attributes that are individually innocuous and scattered across documents. Such a leak
corresponds to no discrete field a filter can excise, which makes it both hard to prevent and,
as this thesis argues, hard to measure.

\section{Problem statement}
\label{sec:intro-problem}

The substantive question is therefore threefold. How much sensitive data does a baseline
enterprise RAG deployment leak, under what attack conditions, and against which kinds of
attacker? How much do existing defense mechanisms actually reduce that leakage when measured
under matched conditions? And --- a question the others quietly depend on --- how do we know the
measurement itself is trustworthy rather than an artifact of one instrument's blind spots? This
last question turns out to be central: a defended RAG system can look almost perfectly protected
or substantially leaky depending only on which instrument does the measuring, and resolving that
ambiguity is as much a part of the contribution as the leakage numbers themselves.

\section{The research gap}
\label{sec:intro-gap}

The literature relevant to this problem divides into five active areas --- RAG-specific
vulnerabilities, PII memorisation and leakage in language models, prompt-injection taxonomies,
defense mechanisms and guardrail architectures, and LLM-as-judge evaluation. Each is mature
within its own boundaries, but, as the review in Chapter~\ref{ch:sota} establishes across fifty
primary sources, none addresses their intersection: where RAG-specific attacks meet
PII-specific privacy metrics, a formal threat taxonomy, layered defense configurations, and
automated leak adjudication. The five cross-cutting gaps that fall out of that synthesis map
one-to-one onto the research questions below.

\section{Research questions}
\label{sec:intro-rqs}

\begin{description}[leftmargin=3.2em, itemsep=2pt]
  \item[\textbf{RQ1}] What privacy leakage occurs in an undefended enterprise RAG baseline, and
    which attack families dominate it? (Chapter~\ref{ch:baseline})
  \item[\textbf{RQ2}] How much do existing defense mechanisms reduce that leakage under matched
    conditions? (Chapter~\ref{ch:defenses})
  \item[\textbf{RQ3}] Can a language-model judge detect and classify privacy leakage reliably,
    with bounded false positives? (Chapter~\ref{ch:defenses})
  \item[\textbf{RQ4}] Which attacker profile is hardest to defend against, and which defense
    works best for each? (Chapter~\ref{ch:tradeoff})
  \item[\textbf{RQ5}] What privacy risk remains after the full defense stack, and how does it
    trade off against utility? (Chapter~\ref{ch:tradeoff})
\end{description}

These are not five independent studies but one linked argument: RQ2 is measured as a reduction
from the RQ1 baseline; RQ5 depends on RQ2 and on utility; and RQ3 underwrites all of them,
because every leakage number is only as trustworthy as the instrument that produced it.

\section{Contributions}
\label{sec:intro-contrib}

The contribution of this thesis is \textbf{convergence, not invention}. It does not propose a
new attack technique or a new defense mechanism; the attacks are drawn from the published
literature and the defense layers are existing, deployed tools. What is missing from the
surveyed work, and what this thesis builds, is the integration of five capabilities into one
framework applied to RAG: a mathematically formalised, LINDDUN-aligned leakage metric pair
$(L,S)$ and a profile-dependent evaluation engine (Chapter~\ref{ch:formal}); a systematic
adversarial campaign across eight attack families (Chapter~\ref{ch:baseline}); a layered defense
whose effect is quantified with $\Delta S$ (Chapters~\ref{ch:defenses}, \ref{ch:tradeoff}); and a
validated language-model judge as the leakage adjudicator (Chapter~\ref{ch:defenses}). No prior
work in the survey couples these for RAG (Chapter~\ref{ch:sota}).

Its sharpest contribution is methodological, and we state it as an \emph{argument} because that
is what it is. We argue that \textbf{measuring a defended RAG system with a single instrument is
wrong in a direction the analyst cannot see}: a deterministic detector, blind to the semantic
and cross-document leakage that survives a defense, records its own blindness as protection.
This thesis demonstrates the effect on a controlled corpus (Chapter~\ref{ch:defenses}) and
argues, in the discussion (Chapter~\ref{ch:discussion}), that the \emph{direction} of the error
generalises beyond it, even though the \emph{magnitude} does not. Three things this thesis is
explicitly \emph{not}: it is not a new attack, not a new defense, and not a certified, deployable
product --- a boundary the findings below make concrete.

\section{Findings, honestly previewed}
\label{sec:intro-findings}

The headline result is a genuine but partial reduction. Through the validated judge --- the
primary instrument --- the full defense stack reduces leakage from $28.8\%$ to $17.3\%$, a
reduction of roughly forty percent, \textbf{not an elimination} (the deterministic engine, by
contrast, reports near-zero, which is precisely the single-instrument error above). The residual
is not uniform: re-identification by \emph{linkage} is the dominant unreduced component ---
redaction defeats direct extraction but cannot remove a person assembled from scattered
attributes (Chapters~\ref{ch:defenses}, \ref{ch:tradeoff}). And the instrument that delivers the
headline is itself bounded: the judge agrees with human annotation at Cohen's $\kappa = 0.827$ on
undefended output but degrades to $\kappa = 0.306$ on the defended distribution the numbers are
built from, so the $17.3\%$ figure is reported with an honest uncertainty interval that contains
it, and the tempting claim that the true rate is \emph{lower} is explicitly refused
(Chapter~\ref{ch:defenses}). The right way to read the whole is as \textbf{a demonstration of a
measurement methodology on a controlled, synthetic corpus, not a certification that any
deployable system is safe}: the relative findings --- which layer matters, which adversary is
hardest, how the residual is composed --- are what transfer; the absolute percentages are not
predictions.

\section{Thesis structure}
\label{sec:intro-structure}

The remainder proceeds as follows. Chapter~\ref{ch:sota} reviews the fifty-source literature
across the five pillars and identifies the gap. Chapter~\ref{ch:threat} establishes the threat
model --- four attacker profiles --- and the formal $(L,S)$ foundations, the synthetic corpus, and
its annotation. Chapter~\ref{ch:baseline} builds the attack framework and measures the
undefended baseline (RQ1). Chapter~\ref{ch:defenses} assembles the layered defense, measures its
effect through two instruments, and validates the judge (RQ2, RQ3). Chapter~\ref{ch:tradeoff}
quantifies the security--utility trade-off and the residual risk (RQ4, RQ5).
Chapter~\ref{ch:discussion} synthesises the five questions into a single argument and states the
limitations, and Chapter~\ref{ch:conclusion} concludes with the contribution and future work.

\ifdefined\ThesisMaster\else
  \bibliographystyle{plain}
  \bibliography{references}
  \end{document}
\fi
